\chapter{\texorpdfstring{\latestrevision{Reproducibility Summary}}{Reproducibility Summary}}
\label{app:reproducibility-summary}
\label{app:experiment-configurations}

{\color{red}
This appendix records the protocol choices needed to interpret and reproduce
the main experimental blocks. It replaces the previous configuration archive
and checkpoint-step registry with a compact description of what was held fixed,
what changed between campaigns, and how checkpoints were selected. Exact
machine-readable configurations and evaluation JSON files remain in the
project repository.

\section{Campaign protocols}

\begin{table}[H]
  \color{red}
  \captionsetup{labelfont={color=red},textfont={color=red}}
  \centering
  \scriptsize
  \caption{Training, selection, and evaluation protocols used by the main
  experimental blocks.}
  \label{tab:reproducibility-campaigns}
  \setlength{\tabcolsep}{3pt}
  \begin{tabular}{L{0.16\textwidth}L{0.23\textwidth}L{0.33\textwidth}L{0.20\textwidth}}
    \toprule
    Campaign & Environment and budget & Checkpoint selection and evaluation & Statistical scope \\
    \midrule
    Stable MLP baseline
      & \texttt{bus14}; 15M steps; 20 parallel environments
      & Periodic deterministic 20-chronic training-split rollouts; retain the
        highest rollout mean
      & 80/20 chronic split; multi-seed core comparisons \\
    \addlinespace
    Preliminary GINE viability
      & \texttt{bus14}; 15M steps; 2{,}000 rollout steps
      & Deterministic evaluation every 80{,}000 steps under the same split as
        the MLP baseline
      & Exploratory architecture study \\
    \addlinespace
    Screens B--F
      & \texttt{bus14}; 15M steps; 72 environments and 576 rollout steps
      & Evaluation every 82{,}944 steps on a changing 20-chronic training
        rollout; selected checkpoint evaluated on 201 held-out chronics
      & One seed per factorial cell (seed 0) \\
    \addlinespace
    Input-distribution study
      & \texttt{bus14} and WCCI; 6{,}000 do-nothing steps; 300-step cap per chronic
      & No learned checkpoint; every visible local-graph node is recorded
      & 426{,}000 and 1{,}236{,}000 node observations, respectively \\
    \addlinespace
    WCCI transfer and adaptation
      & \texttt{bus36\_wcci\_nomaint}; zero-shot, 15M-step scratch, or about
        5M-step fine-tuning
      & Zero-shot uses the selected source checkpoint; WCCI-trained actors use
        their final checkpoint
      & Seed 0; the same 50 chronics, split into 26 easy and 24 difficult cases \\
    \bottomrule
  \end{tabular}
\end{table}

\section{Shared MAPPO settings}

The graph-design screens use the same optimization recipe so that each
factorial changes only the named architectural choice:
\begin{table}[H]
  \color{red}
  \captionsetup{labelfont={color=red},textfont={color=red}}
  \centering
  \small
  \caption{MAPPO settings shared by Screens B--F.}
  \label{tab:reproducibility-mappo}
  \setlength{\tabcolsep}{5pt}
  \begin{tabular}{L{0.27\textwidth}L{0.63\textwidth}}
    \toprule
    Component & Fixed setting \\
    \midrule
    Optimization
      & 10 update epochs, 16 minibatches, learning rate $10^{-4}$ with linear
        decay, maximum gradient norm 1.0 \\
    Returns and PPO
      & $\gamma=0.9$, GAE $\lambda=0.95$, clip ratio 0.2, target KL 0.02,
        entropy coefficient 0.01 \\
    Actor and critic
      & Shared graph encoder with two 128-unit action-head layers; centralized
        three-layer, 256-unit MLP critic \\
    Preprocessing
      & Normalized flat observations and rewards; graph physical scaling and
        sparse-control mechanisms disabled during the screens \\
    Exploration
      & No initial do-nothing prior and no action-0 logit bonus \\
    \bottomrule
  \end{tabular}
\end{table}

\section{Exploratory GINE variants}
\label{app:gine-viability-details}

The preliminary viability study varied several choices simultaneously before
the controlled Screens B--F. Table~\ref{tab:gine-reproducibility-variants}
retains the principal configurations needed to interpret that exploratory
comparison; isolated development variants remain in the machine-readable run
configurations.

\begin{table}[H]
  \color{red}
  \captionsetup{labelfont={color=red},textfont={color=red}}
  \centering
  \scriptsize
  \caption{Principal configurations in the preliminary GINE viability study.}
  \label{tab:gine-reproducibility-variants}
  \label{tab:gine-run-configurations}
  \label{tab:heavy-gine-configurations}
  \setlength{\tabcolsep}{3pt}
  \begin{tabular}{L{0.23\textwidth}L{0.13\textwidth}L{0.12\textwidth}L{0.16\textwidth}L{0.28\textwidth}}
    \toprule
    Run & GINE size & Critic & Actor/critic heads & Main difference from the heavy baseline \\
    \midrule
    \texttt{Heavy GINE Baseline}
      & $2\times128$ & GNN & $3\times256$ / $3\times256$
      & Flat concatenation, edge pre-encoder, entropy 0.02, action-0 prior 0.7,
        legacy critic updates \\
    \texttt{Light GINE (Concat, MLP Critic)}
      & $1\times64$ & MLP & $2\times128$ / $2\times128$
      & Smaller actor, no edge pre-encoder, optimized critic updates \\
    \texttt{Heavy GINE (No Bias)}
      & $2\times128$ & GNN & $3\times256$ / $3\times256$
      & Heavy baseline with the initial action-0 prior removed \\
    \texttt{No Bias, Opt Critic}
      & $2\times128$ & GNN & $3\times256$ / $3\times256$
      & No action-0 prior and optimized critic updates \\
    \texttt{Optimized Heavy GINE}
      & $2\times128$ & GNN & $3\times256$ / $3\times256$
      & No action-0 prior, entropy 0.01, optimized critic updates \\
    \texttt{Optimized Light GINE}
      & $1\times64$ & MLP & $2\times128$ / $2\times128$
      & No flat concatenation or edge pre-encoder; entropy 0.01, no action-0
        prior, optimized critic updates \\
    \bottomrule
  \end{tabular}
\end{table}

Table~\ref{tab:gine-rerun-summary} records the principal outcomes retained
from this exploratory campaign. Figure~\ref{fig:gine-best-search2} gives the
broader training-curve comparison over isolated development variants.

\begin{table}[H]
  \color{red}
  \captionsetup{labelfont={color=red},textfont={color=red}}
  \centering
  \small
  \caption{Survival summary for the principal preliminary GINE variants.}
  \label{tab:gine-rerun-summary}
  \resizebox{\textwidth}{!}{%
  \begin{tabular}{llrr}
    \toprule
    Run & Intended change & Final survival (\%) & Peak survival (\%) \\
    \midrule
    \texttt{Heavy GINE Baseline} & Historical GINE baseline & 59.2 & 93.9 \\
    \texttt{Light GINE (Concat, MLP Critic)} & Light GINE with MLP critic & 78.9 & 90.4 \\
    \texttt{No Bias, Opt Critic} & Optimized critic and bias 0.0 & 82.5 & 85.7 \\
    \texttt{Optimized Light GINE} & Light optimized GINE with MLP critic & 95.7 & 97.4 \\
    \bottomrule
  \end{tabular}
  }
\end{table}

\begin{figure}[H]
  \color{red}
  \captionsetup{labelfont={color=red},textfont={color=red}}
  \centering
  \includegraphics[width=\textwidth]{figures/rerun_gine_best00_vs_best10_14.png}
  \caption{Training comparison between the optimized MLP baseline and the
  strongest shared GINE variants in the preliminary viability study.}
  \label{fig:rerun-gine-best00}
\end{figure}

\begin{figure}[H]
  \color{red}
  \captionsetup{labelfont={color=red},textfont={color=red}}
  \centering
  \includegraphics[width=\textwidth,height=0.82\textheight,keepaspectratio]{figures/gine_best_search2_baseline_comparisons.png}
  \caption{Broader preliminary GINE comparison. Each subplot compares the
  MLP baseline with one exploratory variant. The campaign mixes several
  architectural and optimization changes and is therefore used as supporting
  evidence rather than as a controlled factorial.}
  \label{fig:gine-best-search2}
\end{figure}

\section{Checkpoint and artifact provenance}

For \texttt{bus14}, checkpoint selection uses only the training split. At every
evaluation, a new set of 20 training chronics is rolled out, and the checkpoint
score is their mean survival. The selected checkpoint is then evaluated on the
held-out chronics used by the corresponding experiment. The graph-design
screens use all 201 held-out chronics for this comparison.

Zero-shot WCCI evaluation reuses the selected \texttt{bus14} checkpoint without
target updates. Actors trained or fine-tuned on WCCI are reported at their
final checkpoint, and the transfer comparisons use the same 50 chronics and
seed 0. Exact checkpoint identifiers, per-seed global steps, action summaries,
and evaluation outputs remain in the project artifacts rather than being
duplicated as a thesis table. This also avoids mixing MAPPO environment steps
with behavioral-cloning optimizer steps, which are not directly comparable.
}
